\section{Introdução}\label{sec:intro}

Em um avião monomotor, não há redundância para a propulsão.
Se o motor falha em voo, a aeronave passa a planar e o piloto precisa encontrar onde pousar.
O risco depende da tecnologia do motor.
Em aviões leves na Austrália, o Australian Transport Safety Bureau (ATSB) registrou entre 12 e 13 falhas ou maus funcionamentos \pcmh{} de motores a pistão~\cite{atsb2016}, enquanto os números publicados para turbinas ficam abaixo de 1 \pcmh{}~\cite{stobie2019,cfr25k,easa2017}.
Parte da explicação está na arquitetura: o motor a pistão depende de muitas peças em movimento alternado e de subsistemas em série, e a turbina trabalha em rotação contínua, com mais redundância no controle.

Esses números, porém, são publicados como uma taxa única para o motor inteiro.
Eles não mostram quais subsistemas concentram as falhas, nem permitem avaliar se a redundância, dentro do motor ou na instalação de dois motores, compensa uma tecnologia menos confiável.
Para responder a essas perguntas, é preciso modelar o motor como um sistema.

Este trabalho compara a confiabilidade de motores a pistão, turboélice e turbofan por meio de diagramas de blocos de confiabilidade (RBD, \emph{reliability block diagram}).
Os objetivos específicos, definidos na proposta do projeto, são:
(i)~representar a arquitetura de cada motor em um RBD;
(ii)~estimar a taxa de falha $\lambda$, o tempo médio entre falhas (MTBF) e a confiabilidade $R(t)$ de cada um;
(iii)~avaliar o efeito das redundâncias, isto é, da ignição duplicada no pistão, do controle eletrônico de dois canais no turbofan e do uso de dois motores na mesma aeronave;
e (iv)~identificar os subsistemas que mais contribuem para a falha de cada tecnologia.

Além de responder a esses objetivos, o trabalho mede as frações de cada subsistema em três bases públicas, calcula o valor da causa comum entre motores a partir do qual dois motores a pistão deixam de ser mais seguros que um a turbina e confronta o modelo com a proporção de falhas de motor que terminaram em acidente no Brasil.
